\subsection{Stack Tecnológico, Seguridad y Arquitectura del Sistema}

En esta sección se detallará por completo el stack tecnológico, las medidas de seguridad y la arquitectura que se eligieron para el desarrollo del sistema, junto con una breve justificación e implementación de este. 

\subsubsection{Tecnologías del Backend}

\textbf{TypeScript.} Se escogió este lenguaje de programación para el backend, cuya principal característica es ser un lenguaje de tipado estático ideal para operaciones clave, donde el tipo de valor no debe variar. \par

\textbf{NestJS.} Framework basado en el uso de Node.js orientado a servidores backend para la construcción progresiva de aplicaciones de mediano a gran escalado empresarial, ofreciendo una estructura base para proyectos, además de su arquitectura modular que permitirá separar responsabilidades de manera efectiva. A diferencia del framework de Express, que no tiene una estructura definida inicialmente. \par

\textbf{Fastify.} Se optó por usar el motor HTTP de Fastify para mejor tiempo y serialización de respuesta JSON, reemplazando el motor de Express que NestJS usaba por defecto. \par

\textbf{PrismaORM.} Librería escogida como puente de comunicación entre la base de datos relacional y el servidor backend para la persistencia y lectura de datos, teniendo gran sinergia con TypeScript; esto asegura la integridad de los datos y protección ante inyecciones SQL. \par

\textbf{Argon2.} Algoritmo de hash de contraseñas para la autenticación de los empleados que ingresan a la aplicación mediante su usuario, además de la protección extra de la aplicación frente a ataques mediante el uso de hardware. \par

\textbf{JWT y Passport.} Se utilizaron JWT y Passport para el control de sesiones y acceso de las peticiones a la aplicación, permitiendo que la autenticación no dependa de lecturas y solicitudes constantes a la base de datos. \par

\textbf{BullMQ e IORedis.} Herramientas para la gestión de tareas asincrónicas y procesamiento de tareas pesadas en segundo plano, dando una alta disponibilidad a la aplicación y concurrencia a las peticiones de los usuarios para este tipo de tareas. \par

\textbf{WebSocket.} Protocolo de comunicación usado para el envío de notificaciones de tareas hacia los usuarios, mensajes administrativos o eventos cercanos. \par

\textbf{Librerías de Procesamiento y Reportes (pdfmake, ExcelJS, csv-parser).} Librerías para la lectura y escritura de documentos empresariales para su posterior procesamiento en procesos de alta demanda en el servidor. \par

\textbf{Swagger y OpenAPI.} Para la documentación de la API se usaron estas tecnologías para detallar de manera técnica y automática el contrato de los endpoints establecidos en el sistema. \par

\textbf{Postman.} Herramienta de pruebas y validación de los endpoints, simulando solicitudes HTTP, permitiendo supervisar estados, procesos y respuestas antes de la integración con el frontend. \par

\textbf{Jest.} Framework de pruebas para garantizar la funcionalidad de los procesos y la lógica de negocio, realizando pruebas aisladas que permiten minimizar regresiones en el código y asegurar que las funciones garanticen una respuesta predecible y acorde a su desarrollo. \par

\textbf{Dotenv.} Librería de manejo de variables de entorno para garantizar la seguridad de llaves de servicios, firmas de encriptado y contraseñas de base de datos.

\subsubsection{Tecnologías del Frontend}
\textbf{React 19.} Librería de JavaScript/TypeScript utilizada para la creación de componentes reutilizables y gestionar el renderizado de estos de manera eficiente.\par
\textbf{TypeScript 6.} Lenguaje de tipado estático que permite la integridad y asegurar el tipo de valores de los datos, para su posterior envío al backend.\par
\textbf{Vite 8.} Empaquetador, servidor de desarrollo rápido para la compilación de la aplicación.\par
\textbf{Tailwind CSS.} Framework basado en CSS, que permite diseñar interfaces modernas sin salir del código HTML/TSX.\par
\textbf{Shadcn/ui.} Colección de componentes reutilizables ya creados anteriormente, altamente personalizable (botones, tablas, entre otros).\par
\textbf{lucide-react.} Librería de iconos vectoriales ligeros, que permite dar orientación al usuario dentro de la aplicación.\par
\textbf{TanStack Query 5.} Sincronización de caché, carga y actualización de datos que vienen por parte del servidor.\par
\textbf{Zustand 5.} Librería de gestión de estado global, que permite guardar sesión y compartir datos entre pantallas.\par
\textbf{Axios.} Cliente HTTP para realizar solicitudes al backend y obtener respuestas de manera eficiente.\par
\textbf{React Hook Form 7.} Librería para el manejo de formularios, reduciendo cargas innecesarias al escribir.\par
\textbf{Zod.} Librería de validación de datos, permite la comparación y estructura del envío de datos al backend.\par
\textbf{React Router DOM 7.} Sistema de navegación para gestionar distintas páginas de la aplicación web sin recargar la página completa.\par
\subsubsection{Paquete de validación de Schemas compartido.} 
Paquete independiente que permite centralizar la validación de schemas, tipos de datos y reglas de estas. \par

\textbf{ZOD.} Librería para validación de datos, perfecta para establecer contratos de entrada a los endpoints de la aplicación.

\subsubsection{Bases de Datos}
Se escogió usar una base de datos para permanencia de registros historicos, centralizacion, lectura y auditoria de estos.

\textbf{PostgreSQL.} Motor de base de datos relacional elegido por su robustez, ideal para proyectos empresariales garantizando la estructura de guardado de los datos de los usuarios. \par

\textbf{Redis.} Sistema de almacenamiento de datos en memoria, de alto rendimiento en procedimientos de caché y colas, para la concurrencia en el servidor. \par

\textbf{UUIDv7.} Método de identificación a entidades dentro de la base de datos relacional; a diferencia de sus versiones anteriores que generan el identificador de manera aleatoria, este realiza la generación mediante orden cronológico, lo cual ayuda a auditar la creación de registros.

\subsubsection{Criptografía y Seguridad de Datos}
En esta sección se explicarán los mecanismos internos para garantizar la integridad de los datos del usuario al momento de hacer persistencia y peticiones dentro y fuera de la base de datos.

\textbf{Enmascaramiento de datos.} Técnica para la protección de datos de los usuarios en las respuestas JSON de los endpoints, que consiste en que una parte del valor que almacena la variable sea reemplazada por otro carácter. \par

\textbf{Cifrado de datos.} Técnica para la protección de datos, consiste en transformar un dato legible a uno ilegible antes de ser enviado a la base de datos para garantizar que, en caso de una fuga de información, los datos no sean capaces de ser leídos por el atacante.

\subsubsection{Diagrama de Arquitectura de Software}
En el siguiente diagrama se observará la estructura y comportamiento general del sistema, la estructura del proyecto se basa en la arquitectura de un monolito modular combinado con Clean Architecture, donde los controladores, mediante el motor de peticiones de Fastify, procesan y validan las peticiones mediante JWT, además de las tareas de alta demanda (cálculo de planillas) se desarrollan mediante colas usando BullMQ, donde finalmente se designa a PrismaORM para la comunicación con la base de datos que es PostgreSQL.
\vspace{0.3cm}
\par\medskip
\noindent
\figuraEnTexto
    {fig:arquitectura_backend}
    {Diagrama de arquitectura general y flujo del backend}
    {images/diagramaFuncionamientoGeneral.png}
    {Estructura modular y flujo de comunicación del backend. Elaboración propia.}
\par\medskip
